\input{../tex-inputs/latex-korrekturansicht-vorspann}

\section[Berta Zuckerkandl an Arthur Schnitzler, {[}zwischen 11. und 31. 5. 1911?{]}]{L03997 Berta Zuckerkandl an Arthur Schnitzler, {[}zwischen 11. und 31. 5. 1911?{]}}
\nopagebreak\mylabel{L03997v}
\rehead{ }\normalsize\beginnumbering\briefempfaengerindex{Schnitzler, Arthur@\textsc{Schnitzler, Arthur}!zzzZuckerkandl, Berta@\emph{von Berta Zuckerkandl}!1911-05-311@{{[}zwischen 11. und 31. 5. 1911?{]}}|(be}
\toendnotes[C]{\smallbreak\pagebreak[2]}
\correspDesc{Versand  durch Berta Zuckerkandl im Zeitraum [zwischen 11. und
                  31. 5. 1911?] in Wien
\newline{}Erhalt  durch Arthur Schnitzler in Wien}\toendnotes[C]{\smallbreak}
\Standort{CUL, Schnitzler, B 200.}
\physDesc{Brief, 1 Blatt, 4 Seiten, 721 Zeichen
\newline{}Handschrift: schwarze Tinte, lateinische Kurrent
\newline{}Schnitzler: mit Bleistift Vermerk: »\textsc{Zuckerkandl}.« }\toendnotes[C]{\smallbreak}
\pstart{}{\pb}Hochverehrter Herr Doktor!\pend\vspace{0.5em}
\pstart
           Ich habe \textcolor{green}{es}\pwindex{Schnitzler, Arthur 15. 5. 1862 Wien – 21. 10. 1931 ebd.@\textsc{Schnitzler, Arthur} (15. 5. 1862 Wien – 21. 10. 1931 ebd.), \emph{Schriftsteller, Mediziner}!weite Land. Tragikomödie in fünf Akten@\strich\emph{Das weite Land. Tragikomödie in fünf Akten}|pwv}{}\ledrightnote{{$\rightarrow$}\emph{\textcolor{green}{Das weite Land. Tragikomödie in fünf Akten}}} erst in einem Zug
               gelesen. Und dann langsam nachgekostet. Es ist ganz wunderbar. Wirklich \label{K_L03997-1v}\edtext{ein weites – \textcolor{green}{weites Land}\pwindex{Schnitzler, Arthur 15. 5. 1862 Wien – 21. 10. 1931 ebd.@\textsc{Schnitzler, Arthur} (15. 5. 1862 Wien – 21. 10. 1931 ebd.), \emph{Schriftsteller, Mediziner}!weite Land. Tragikomödie in fünf Akten@\strich\emph{Das weite Land. Tragikomödie in fünf Akten}|pwv}{}\ledrightnote{{$\rightarrow$}\emph{\textcolor{green}{Das weite Land. Tragikomödie in fünf Akten}}}}{\lemma{\textnormal{\emph{ein weites – weites Land}}}\Cendnote{\textnormal{Die Druckausgabe von \emph{\textcolor{green}{Das weite Land}\pwindex{Schnitzler, Arthur 15. 5. 1862 Wien – 21. 10. 1931 ebd.@\textsc{Schnitzler, Arthur} (15. 5. 1862 Wien – 21. 10. 1931 ebd.), \emph{Schriftsteller, Mediziner}!weite Land. Tragikomödie in fünf Akten@\strich\emph{Das weite Land. Tragikomödie in fünf Akten}|pwk}} erschien erst im September 1911, es muss sich hier also um das Bühnenmanuskript handeln.  (\textcolor{blue}{Schnitzler} hatte es auch \textcolor{blue}{Stefan Zweig}\pwindex{Zweig, Stefan 28.\,11.\,1881 Wien – 23.\,2.\,1942 Petrópolis@\textsc{Zweig, Stefan} (28.\,11.\,1881 Wien – 23.\,2.\,1942 Petrópolis), \emph{Schriftsteller}|pwk} bereits zur Verfügung gestellt, vgl. Stefan Zweig an Arthur Schnitzler, [5. oder 6. 10. 1910?].)}}}\label{K_L03997-1} in das wir
               durch Dichter’s Gnaden {\pb}blicken können.
               Wie viel wird lebendig in der eigenen Seele!\pend
           
\pstart
           – Ich glaube dass gerade dieses Thema – das Verhältniss zwischen Mann und Frau – in
               der so eigenthümlichen {\pb}Beleuchtung – für
                  \textcolor{pink}{Paris}\oindex{Paris@\textbf{Paris}, \emph{Hauptstadt}|pw}{}\ledrightnote{\textcolor{pink}{Paris}} wie geschaffen wäre. Erlauben Sie mir
               jedenfalls wenn ich hinreise – dort darüber unverbindlich zu sprechen. Übrigens
               erwarte ich täglich einen Brief – \label{K_L03997-2v}\edtext{wegen des \textcolor{green}{Medardus}\pwindex{Schnitzler, Arthur 15. 5. 1862 Wien – 21. 10. 1931 ebd.@\textsc{Schnitzler, Arthur} (15. 5. 1862 Wien – 21. 10. 1931 ebd.), \emph{Schriftsteller, Mediziner}!junge Medardus. Dramatische Historie in einem Vorspiel und fünf Aufzügen@\strich\emph{Der junge Medardus. Dramatische Historie in einem Vorspiel und fünf Aufzügen}|pw}{}\ledrightnote{\textcolor{green}{Der junge Medardus. Dramatische Historie in einem Vorspiel und fünf Aufzügen}}}{\lemma{\textnormal{\emph{wegen des Medardus}}}\Cendnote{\textnormal{Im April 1911 war \textcolor{blue}{Berta Zuckerkandl}\pwindex{Zuckerkandl, Berta 13.\,4.\,1864 Wien – 16.\,10.\,1945 Paris@\textsc{Zuckerkandl, Berta} (13.\,4.\,1864 Wien – 16.\,10.\,1945 Paris), \emph{Schriftstellerin, Journalistin, Übersetzerin}|pwk} mit dem Vorschlag an \textcolor{blue}{Schnitzler} herangetreten, sein Schauspiel \emph{\textcolor{green}{Der junge Medardus}\pwindex{Schnitzler, Arthur 15. 5. 1862 Wien – 21. 10. 1931 ebd.@\textsc{Schnitzler, Arthur} (15. 5. 1862 Wien – 21. 10. 1931 ebd.), \emph{Schriftsteller, Mediziner}!junge Medardus. Dramatische Historie in einem Vorspiel und fünf Aufzügen@\strich\emph{Der junge Medardus. Dramatische Historie in einem Vorspiel und fünf Aufzügen}|pwk}} zur Aufführung in \textcolor{pink}{Paris}\oindex{Paris@\textbf{Paris}, \emph{Hauptstadt}|pwk} vermitteln (vgl. Berta Zuckerkandl an Arthur Schnitzler, [14. 4. 1911?]), am 25. 7. 1911 unterrichtete sie ihn, dass der
                  Versuch vorerst gescheitert war (vgl. Berta Zuckerkandl an Arthur Schnitzler, 25. 7. [1911]). Der vorliegende undatierte Brief, in dem \textcolor{blue}{Zuckerkandl}\pwindex{Zuckerkandl, Berta 13.\,4.\,1864 Wien – 16.\,10.\,1945 Paris@\textsc{Zuckerkandl, Berta} (13.\,4.\,1864 Wien – 16.\,10.\,1945 Paris), \emph{Schriftstellerin, Journalistin, Übersetzerin}|pwk} angibt auf Nachrichten bezüglich des \textcolor{green}{Medardus}\pwindex{Schnitzler, Arthur 15. 5. 1862 Wien – 21. 10. 1931 ebd.@\textsc{Schnitzler, Arthur} (15. 5. 1862 Wien – 21. 10. 1931 ebd.), \emph{Schriftsteller, Mediziner}!junge Medardus. Dramatische Historie in einem Vorspiel und fünf Aufzügen@\strich\emph{Der junge Medardus. Dramatische Historie in einem Vorspiel und fünf Aufzügen}|pwkv} zu warten, fällt
                  also in diese Zeitspanne. Genauer: Am 18. 4. 1911 schlug \textcolor{blue}{Schnitzler} vor, Anfang Mai über
                  die Vermittlung von \emph{\textcolor{green}{Der junge Medardus}\pwindex{Schnitzler, Arthur 15. 5. 1862 Wien – 21. 10. 1931 ebd.@\textsc{Schnitzler, Arthur} (15. 5. 1862 Wien – 21. 10. 1931 ebd.), \emph{Schriftsteller, Mediziner}!junge Medardus. Dramatische Historie in einem Vorspiel und fünf Aufzügen@\strich\emph{Der junge Medardus. Dramatische Historie in einem Vorspiel und fünf Aufzügen}|pwk}} aber
                  auch von seinem neuen Stück \emph{\textcolor{green}{Das weite Land}\pwindex{Schnitzler, Arthur 15. 5. 1862 Wien – 21. 10. 1931 ebd.@\textsc{Schnitzler, Arthur} (15. 5. 1862 Wien – 21. 10. 1931 ebd.), \emph{Schriftsteller, Mediziner}!weite Land. Tragikomödie in fünf Akten@\strich\emph{Das weite Land. Tragikomödie in fünf Akten}|pwk}}
                  nach \textcolor{pink}{Frankreich}\oindex{Frankreich@\textbf{Frankreich}|pwk} zu sprechen: »Dieses
                        \textcolor{green}{Stück}\pwindex{Schnitzler, Arthur 15. 5. 1862 Wien – 21. 10. 1931 ebd.@\textsc{Schnitzler, Arthur} (15. 5. 1862 Wien – 21. 10. 1931 ebd.), \emph{Schriftsteller, Mediziner}!weite Land. Tragikomödie in fünf Akten@\strich\emph{Das weite Land. Tragikomödie in fünf Akten}|pwv} scheint mir
                     nach den internationalen Seite mehr zu versprechen als der \textcolor{green}{Medardus}\pwindex{Schnitzler, Arthur 15. 5. 1862 Wien – 21. 10. 1931 ebd.@\textsc{Schnitzler, Arthur} (15. 5. 1862 Wien – 21. 10. 1931 ebd.), \emph{Schriftsteller, Mediziner}!junge Medardus. Dramatische Historie in einem Vorspiel und fünf Aufzügen@\strich\emph{Der junge Medardus. Dramatische Historie in einem Vorspiel und fünf Aufzügen}|pw}«, siehe Arthur Schnitzler an Berta Zuckerkandl, 18. 4. 1911. Das
                  Treffen fand am 11. 5. 1911 statt. Der vorliegende Brief dürfte nach dem Treffen
                  entstanden sein (möglicherweise gab \textcolor{blue}{Schnitzler}{ }\textcolor{blue}{Zuckerkandl}\pwindex{Zuckerkandl, Berta 13.\,4.\,1864 Wien – 16.\,10.\,1945 Paris@\textsc{Zuckerkandl, Berta} (13.\,4.\,1864 Wien – 16.\,10.\,1945 Paris), \emph{Schriftstellerin, Journalistin, Übersetzerin}|pwk} den \textcolor{green}{Text}\pwindex{Schnitzler, Arthur 15. 5. 1862 Wien – 21. 10. 1931 ebd.@\textsc{Schnitzler, Arthur} (15. 5. 1862 Wien – 21. 10. 1931 ebd.), \emph{Schriftsteller, Mediziner}!weite Land. Tragikomödie in fünf Akten@\strich\emph{Das weite Land. Tragikomödie in fünf Akten}|pwkv} gleich mit) und vor \textcolor{blue}{Zuckerkandls}\pwindex{Zuckerkandl, Berta 13.\,4.\,1864 Wien – 16.\,10.\,1945 Paris@\textsc{Zuckerkandl, Berta} (13.\,4.\,1864 Wien – 16.\,10.\,1945 Paris), \emph{Schriftstellerin, Journalistin, Übersetzerin}|pwk} erstem Brief aus dem
                     Juni, in dem sie berichtet Nachricht aus \textcolor{pink}{Paris}\oindex{Paris@\textbf{Paris}, \emph{Hauptstadt}|pwk} erhalten zu haben, und zwar mindestens 10 Tage vor
                  diesem Brief, denn dort entschuldigt sie sich für die lange Pause: »Ich war
                     sehr leidend und arbeitsunfähig. Auch hatte ich nichts Weiteres aus \textcolor{pink}{Paris}\oindex{Paris@\textbf{Paris}, \emph{Hauptstadt}|pw} gehört.« (siehe Berta Zuckerkandl an Arthur Schnitzler, [zwischen 11. und
               13. 6. 1911?]). So ergibt sich die
                  ungefähre Datierung zwischen 11. und 31. 5. 1911.}}}\label{K_L03997-2}.\pend
           
\pstart
           {\pb}Man sagt i{\geminationm}er Frauen seien unergründlich. Ihre Männertypen – schillern noch rätselhafter.
                  \label{K_L03997-3v}\edtext{\textcolor{green}{Friedrich}\pwindex{Schnitzler, Arthur 15. 5. 1862 Wien – 21. 10. 1931 ebd.@\textsc{Schnitzler, Arthur} (15. 5. 1862 Wien – 21. 10. 1931 ebd.), \emph{Schriftsteller, Mediziner}!weite Land. Tragikomödie in fünf Akten@\strich\emph{Das weite Land. Tragikomödie in fünf Akten}|pwv}{}\ledrightnote{{$\rightarrow$}\emph{\textcolor{green}{Das weite Land. Tragikomödie in fünf Akten}}}}{\lemma{\textnormal{\emph{Friedrich}}}\Cendnote{\textnormal{Hauptfigur von \emph{\textcolor{green}{Das weite Land}\pwindex{Schnitzler, Arthur 15. 5. 1862 Wien – 21. 10. 1931 ebd.@\textsc{Schnitzler, Arthur} (15. 5. 1862 Wien – 21. 10. 1931 ebd.), \emph{Schriftsteller, Mediziner}!weite Land. Tragikomödie in fünf Akten@\strich\emph{Das weite Land. Tragikomödie in fünf Akten}|pwk}}}}}\label{K_L03997-3} – ! Alles giebt so viel Denck- und Gefühls-Nahrung.\pend
           
\pstart
           Viel herzlich Danck, {\\[\baselineskip]}\spacefill\mbox{B. Zuckerkandl}\pend
           \leftskip=0em{}\selectlanguage{ngerman}\endnumbering\briefempfaengerindex{Schnitzler, Arthur@\textsc{Schnitzler, Arthur}!zzzZuckerkandl, Berta@\emph{von Berta Zuckerkandl}!1911-05-111@{{[}zwischen 11. und 31. 5. 1911?{]}}|)be}\mylabel{L03997h}  \newcommand{\dateiname}{L03997}\newcommand{\titel}{Berta Zuckerkandl an Arthur Schnitzler, [zwischen 11. und 31. 5. 1911?]}\newcommand{\editorInnen}{Herausgegeben von Jahnke, SelmaMüller, Martin Anton}\input{../tex-inputs/latex-korrekturansicht-abspann}
